% slides/secciones/slides_contenido.tex

% --- INTRODUCCIÓN Y CONTEXTO ---
\section{Introducción}
\begin{frame}{Hacia una Nueva Logística Sanitaria}
    \begin{exampleblock}{El Cambio de Paradigma}
        La automatización en hospitales ha evolucionado de una visión futurista a una necesidad operativa para optimizar tiempos y reducir la carga no asistencial.
    \end{exampleblock}
    
    \vspace{0.5cm}
    \textbf{Foco de Investigación:}
    Desarrollo de cobots móviles bajo el estándar ROS 2, priorizando la seguridad biológica y la estabilidad en el transporte de materiales críticos.
\end{frame}

\begin{frame}{Motivación: El Desafío Socio-Sanitario}
    \begin{columns}
        \begin{column}{0.5\textwidth}
            \begin{block}{Entorno Crítico}
                Hospitales saturados con alta demanda de logística interna (medicinas, muestras, residuos).
            \end{block}
        \end{column}
        \begin{column}{0.5\textwidth}
            \begin{block}{Objetivo Humano}
                Liberar al personal de enfermería de tareas de transporte para centrarse en el cuidado directo del paciente.
            \end{block}
        \end{column}
    \end{columns}
    \vspace{0.4cm}
    \textit{Un robot logístico eficiente reduce en un 30\% el tiempo de desplazamiento del personal sanitario.}
\end{frame}

% --- ESTADO DEL ARTE ---
\section{Estado del Arte}
\begin{frame}{Referentes Comerciales Actuales}
    \begin{itemize}
        \item \textbf{Aethon TUG:} Especialista en carga pesada (450kg) e integración con ascensores.
        \item \textbf{MiR (Mobile Industrial Robots):} Flexibilidad y facilidad de configuración en entornos industriales y médicos.
        \item \textbf{Diligent Moxi:} Robótica asistencial con brazo mecánico para interacción con el entorno.
    \end{itemize}
    \vspace{0.3cm}
    \begin{block}{Limitación Común}
        Sistemas rígidos que a menudo requieren infraestructuras muy específicas o carecen de estabilidad para muestras biológicas líquidas.
    \end{block}
\end{frame}

\begin{frame}{Ecosistemas de Simulación y Middleware}
    \textbf{ROS 2 como Estándar Tecnológico:}
    \begin{itemize}
        \item[$\bullet$] \textbf{Nav2:} El stack de navegación más avanzado para evitar obstáculos dinámicos.
        \item[$\bullet$] \textbf{Open-RMF:} Coordinación de flotas multimarca y gestión de recursos compartidos.
        \item[$\bullet$] \textbf{Gazebo / Ignition:} Simulación física de alta fidelidad.
    \end{itemize}
    \vspace{0.3cm}
    \textit{Validar en virtual ahorra meses de desarrollo y garantiza la seguridad antes del despliegue real.}
\end{frame}

% --- DISEÑO DEL ROBOT ---
\section{Propuesta Técnica}
\begin{frame}{HospitalBot v1.0: Nuestra Propuesta}
    \begin{center}
        \textbf{Un robot diseñado por y para el hospital.}
    \end{center}
    \begin{columns}
        \begin{column}{0.5\textwidth}
            \textbf{Concepto:}
            \begin{itemize}
                \item Estética no amenazante.
                \item Cuerpo cilíndrico.
                \item Fácil desinfección.
            \end{itemize}
        \end{column}
        \begin{column}{0.5\textwidth}
            \textbf{Innovación:}
            \begin{itemize}
                \item Suspensión activa.
                \item Materiales grado médico.
                \item Doble sistema de visión.
            \end{itemize}
        \end{column}
    \end{columns}
\end{frame}

\begin{frame}{Diseño Industrial: Materiales y Asepsia}
    \begin{block}{Superficies Antimicrobianas}
        Uso de polímeros técnicos con aditivos de iones de plata y acero inoxidable 316L.
    \end{block}
    \begin{itemize}
        \item[$\checkmark$] Resistencia al peróxido de hidrógeno y alcoholes.
        \item[$\checkmark$] Acabados redondeados para evitar nichos bacterianos.
        \item[$\checkmark$] Compartimentos sellados para evitar derrames biológicos.
    \end{itemize}
\end{frame}

\begin{frame}{Ingeniería Mecánica: Suspensión Independiente}
    \textbf{El reto de las muestras biológicas:}
    Las vibraciones en juntas de dilatación y suelos irregulares pueden degradar muestras de sangre o tejidos.
    
    \vspace{0.3cm}
    \begin{exampleblock}{Solución Técnica}
        Suspensión de doble horquilla con amortiguación variable controlada por software para mantener la carga útil siempre nivelada y libre de vibraciones de alta frecuencia.
    \end{exampleblock}
\end{frame}

\begin{frame}{Sistemas de Percepción: Fusión Sensorial}
    \begin{columns}
        \begin{column}{0.6\textwidth}
            \textbf{Sensores Principales:}
            \begin{itemize}
                \item \textbf{LiDAR 2D:} Mapeo plano y SLAM.
                \item \textbf{Cámara RGB-D:} Visión estereoscópica para detectar objetos bajos (pies, cables).
                \item \textbf{IMU 9-Ejes:} Estabilización y corrección de ruta.
            \end{itemize}
        \end{column}
        \begin{column}{0.4\textwidth}
            \begin{alertblock}{Seguridad}
                Sensores de ultrasonido como sistema de frenado de emergencia independiente del procesador principal.
            \end{alertblock}
        \end{column}
    \end{columns}
\end{frame}

\begin{frame}{Arquitectura de Software: Ecosistema ROS 2}
    \begin{block}{Comunicaciones DDS}
        Uso de \textit{Data Distribution Service} para garantizar latencias bajas y fiabilidad en el intercambio de datos entre nodos de control.
    \end{block}
    \vspace{0.3cm}
    \begin{itemize}
        \item[1.] \textbf{Percepción:} Procesamiento de nubes de puntos.
        \item[2.] \textbf{Planificación:} Algoritmos A* y DWA para rutas fluidas.
        \item[3.] \textbf{Control:} Lazo cerrado con encoders magnéticos.
    \end{itemize}
\end{frame}

% --- SEGURIDAD E INFRAESTRUCTURA ---
\section{Seguridad e Integración}
\begin{frame}{Marco Normativo: ISO 13482:2014}
    \textbf{Requisitos de Seguridad para Robots de Cuidado Personal:}
    \begin{itemize}
        \item[$\bullet$] \textbf{Control de Energía:} Límites estrictos de velocidad en pasillos (máx 1.2 m/s).
        \item[$\bullet$] \textbf{Detección de Personas:} Prioridad absoluta al peatón/paciente.
        \item[$\bullet$] \textbf{Parada de Seguridad:} Hardware redundante certificado.
    \end{itemize}
    \vspace{0.3cm}
    \textit{El robot no solo debe evitar colisiones, debe ser predecible para el humano.}
\end{frame}

\begin{frame}{Infraestructura: El Hospital Inteligente}
    El robot es un "habitante" más del edificio digital.
    \begin{itemize}
        \item \textbf{Gestión de Ascensores:} Protocolos inalámbricos para cambio de planta.
        \item \textbf{Puertas Automáticas:} Apertura remota mediante balizas BLE o WiFi.
        \item \textbf{Docking Station:} Recarga automática por inducción para evitar conectores expuestos (limpieza fácil).
    \end{itemize}
\end{frame}

\begin{frame}{Protocolos Sanitarios: Desinfección UV-C}
    \begin{block}{Mitigación de Infecciones Nosocomiales}
        El robot no puede ser un vector de contagio.
    \end{block}
    \vspace{0.3cm}
    \begin{columns}
        \begin{column}{0.5\textwidth}
            \textbf{Modo Activo:}
            Lámparas UV-C integradas para desinfectar pasillos en horas nocturnas.
        \end{column}
        \begin{column}{0.5\textwidth}
            \textbf{Modo Pasivo:}
            Ciclos de limpieza en salas estancas tras transportar residuos.
        \end{column}
    \end{columns}
\end{frame}

% --- CONCLUSIONES ---
\section{Cierre}
\begin{frame}{Conclusiones y Perspectivas Futuras}
    \begin{exampleblock}{Sinergia Humano-Robot}
        La tecnología actúa como un facilitador que libera al personal sanitario para tareas de alto valor.
    \end{exampleblock}
    
    \vspace{0.4cm}
    \textbf{Mapa de Ruta:}
    \begin{itemize}
        \item[$\rightarrow$] Integración de 5G para telemetría masiva.
        \item[$\rightarrow$] Algoritmos de IA para predicción de flujos de pacientes.
        \item[$\rightarrow$] Despliegue de flotas heterogéneas (Logística + Asistencia).
    \end{itemize}
\end{frame}
